\section{Knowledge organization and lifecycle management}
\label{sec:organization}

Organization determines which candidates retrieval can see and whether the
stored state remains trustworthy. QDCVR places documents by content rather
than filename; split parts inherit the source assignment, and the catalog
exposes library and document descriptions before full bodies are read.

\subsection{Metadata, tags and duplicate review}
\label{sec:org-placement}
\label{sec:org-registries}
\label{sec:org-tags}
\label{sec:org-dedup}

Integrity checks compare the document files, the \texttt{.tree-fs.json} tree
and each library's \texttt{.knowledge-base.yml} for document existence, path
and index metadata. Tags derive from content and are checked against it.
Duplicate detection reports candidate pairs for review; deletion is not
automatic because adjacent parts may be similar while containing distinct
source sections.

Sub-libraries are first-class tree nodes with their own descriptions and
indices. Index writes resolve the document's owning library; searching a parent
library traverses its descendants' collections. This supports hierarchical
organization without implying that child vectors are stored in the parent's
collection.

\subsection{Lifecycle operations, graph and operational knowledge}
\label{sec:org-lifecycle}
\label{sec:org-graph}
\label{sec:org-experience}

Re-indexing is explicit: a rename can leave stale vector chunks until a forced
re-index, so operation completion requires checking the resulting index, not
only the metadata record. The document graph supports navigation and candidate
expansion, not evidence. Operational lessons are stored as retrievable
documents with credibility tiers, but the available measurements are case
examples rather than a scale evaluation of distillation, approval or
credibility decay.

\subsection{Organization contract}
\label{sec:org-integrity}
\label{sec:org-summary}

The evaluation checks three-source consistency, content-derived descriptions,
regression probes and lifecycle outcomes (Section~\ref{sec:e2}). The contract
makes stale or inconsistent state detectable; it does not prevent every
failure by construction. Merge and split remain outside the controlled test
scope, so performance claims must be read within the corpus sizes reported in
the Evaluation section.
